% فصل ۱ — مقدمه
% همهٔ عددها از docs/generated/thesis_numbers.json می‌آیند، که src/thesis_numbers.py آن را
% از آرشیو runs/ و docs/generated/ می‌سازد. صحت با «python3 tools/farsi_lint.py --numbers» بررسی می‌شود.

\chapter{مقدمه}
\label{ch:intro}

\section{مسئله}

رتینوپاتی دیابتی بزرگ‌ترین علت نابینایی پیشگیری‌پذیر در جمعیت شاغل جهان است، و ادم ماکولای
دیابتی شایع‌ترین سازوکاری است که در آن بینایی مرکزی از دست می‌رود. هر دو با غربالگری منظم و
درمان به‌موقع قابل کنترل‌اند، و هر دو تا مرحله‌ای پیش می‌روند که بیمار علامتی حس نکند. پس
مسئله، پیش از آن‌که مسئلهٔ درمان باشد، مسئلهٔ \emph{ظرفیت غربالگری} است: شمار افراد دیابتی از
شمار متخصصانی که می‌توانند ته‌چشمشان را درجه‌بندی کنند بسیار بیشتر است.

همین شکاف، دلیل سه دههٔ کار روی درجه‌بندی خودکار تصویر ته‌چشم است. مدل‌های یادگیری عمیق روی
این وظیفه اعدادی گزارش می‌کنند که در نگاه نخست مسئله را حل‌شده نشان می‌دهد. این پایان‌نامه از
همان‌جا آغاز می‌کند که آن اعداد چه چیزی را نشان می‌دهند و چه چیزی را نشان نمی‌دهند.

\section{شکافی که بسته نمی‌شود}

اگر مسئله حل شده بود، سامانه‌های خودکار در بالین می‌بودند. نیستند. تا پایان آوریل \num{۲۰۲۴} تنها
\textbf{\num{۳}} دستگاه هوش مصنوعی چشم‌پزشکی از سازمان غذا و داروی ایالات متحده مجوز گرفته
بودند، و هر سه تنها برای غربالگری رتینوپاتی دیابتی و تنها با دوربین‌های مشخص
\cite{ruamviboonsuk2024review}. مرور یادشده نتیجه را صریح می‌نویسد: شمار پژوهش‌های روش‌های
تازه نمایی بالا می‌رود، در حالی که پذیرش در مراقبت چشمی با شیبی بسیار کمتر رشد می‌کند، پس
\textbf{شکاف میان توسعه و استقرار در حال گشاد شدن است}.

همان مرور چهار مانع را نام می‌برد، و مانع نخست انگیزهٔ این پایان‌نامه است: \textbf{کمبود
مقایسه‌های سربه‌سر میان مدل‌های موجود}. سه مانع دیگر — نبود شواهد روشن هزینه‌اثربخشی، مسائل
عدالت و سوگیری، و مسئولیت حقوقی-پزشکی — بیرون از دامنهٔ این کارند و در فصل \ref{ch:discussion}
به‌عنوان محدودیت نام برده می‌شوند.

مرور نظام‌مند دیگری، روی \textbf{\num{۳۸}} پژوهش، نشان می‌دهد این کمبود از کجا می‌آید: تنها
\textbf{\num{۱۲}} پژوهش، یعنی \textbf{\num{۳۲}} درصد، اعتبارسنجی بیرونی انجام داده‌اند، و تنها
\textbf{\num{۴}} پژوهش، یعنی \textbf{\num{۱۱}} درصد، شیوهٔ برخورد با داده‌های ازدست‌رفته را گزارش
کرده‌اند \cite{yang2025review}. همان مرور، در بخش بحث، نتیجه‌ای می‌گیرد که در سراسر این
پایان‌نامه تکرار خواهد شد: مدل‌هایی که تنها روی دادهٔ درونی آموزش دیده‌اند کارایی بالاتری
گزارش می‌کنند، که نشانهٔ خطر بیش‌برازش و سوگیری آموزش است.

\section{آنچه یک عدد سرتیتر پنهان می‌کند}

بین «مدل خوب کار می‌کند» و «مدل در بالین کار می‌کند» یک لایهٔ نادیده وجود دارد که این
پایان‌نامه دربارهٔ آن است: \textbf{قاعدهٔ تصمیم}.

یک مدل درجه‌بندی، نمره‌ای پیوسته تولید می‌کند. آنچه بالین می‌بیند یک درجه است، و میان آن دو،
مجموعه‌ای از نقاط برش می‌نشیند. این نقاط تقریباً همیشه انتخاب نمی‌شوند: آستانهٔ پیش‌فرض
\wrongnum{۵/۰} روی خروجی سیگموئید — که در این متن به‌درستی \num{۰٫۵} نوشته می‌شود — از مقیاس
درجه‌بندی نمی‌آید، از تابع فعال‌سازی می‌آید. \textbf{یک پیش‌فرض تنظیم‌نشده، فراپارامتری است
که انتخاب شده، نه فراپارامتری که از آن پرهیز شده است.}

اهمیت این نکته را ادبیات موجود، بی‌آن‌که نتیجه‌اش را بگیرد، بارها نشان داده است. در یک
اعتبارسنجی واقعی از سامانه‌های تجاری، ویژگی سه فروشنده روی \emph{همان چشم‌ها} و با همان مرجع،
از \textbf{\num{۱۴٫۲۵}} درصد تا \textbf{\num{۹۶٫۰۱}} درصد نوسان می‌کند \cite{duggal2025realworld}. در یک
فراتحلیل، حساسیت تجمیعی \textbf{\num{۰٫۹۵}} گزارش می‌شود در حالی که پژوهش‌های واردشده بازه‌ای تا
\textbf{\num{۰٫۰۶}} را در بر می‌گیرند \cite{tahir2025meta}. و در یک محک ترتیبی، همان قاعدهٔ آستانه
که کم‌برآوردی درجه‌های شدید را از \textbf{\num{۸۴٫۷}} به \textbf{\num{۴۹٫۷}} می‌رساند، شمار
پیش‌بینی‌های درجهٔ یک را به \textbf{\num{۷۱۲٫۷}} در برابر \textbf{\num{۲۷۰}} مورد واقعی می‌رساند، با
دقت \textbf{\num{۰٫۱۹۳}} \cite{sheng2026orderdr}.

\textbf{هیچ‌یک از این مقاله‌ها یک جاروب آستانه گزارش نمی‌کند.} همه اثر آن را می‌بینند و
هیچ‌کدام آن را به‌عنوان متغیر مطالعه نمی‌کنند.

\section{پرسش‌های این پژوهش}

از این‌جا سه پرسش می‌آید که ساختار پایان‌نامه را می‌سازند.

\textbf{پرسش یکم — قاعدهٔ تصمیم چقدر از آنچه به معماری نسبت می‌دهیم را توضیح می‌دهد؟} اگر
جابه‌جایی نقاط برش بتواند اختلاف میان دو معماری را از میان ببرد یا علامتش را عوض کند، آنگاه
هر مقایسهٔ معماری که نقاط برش را همتا نکرده باشد، دربارهٔ آستانه حرف می‌زند و فکر می‌کند
دربارهٔ بازنمایی حرف می‌زند.

\textbf{پرسش دوم — سقف کارایی سر ادم ماکولا کجاست، و از داده می‌آید یا از معماری؟} این پرسش
به این دلیل قابل پاسخ است که می‌توان مداخله‌هایی طراحی کرد که مستقیماً ناحیهٔ تعریف‌کنندهٔ
برچسب را هدف بگیرند؛ اگر آن‌ها هم کاری نکنند، محدودیت جای دیگری است.

\textbf{پرسش سوم — تبار داده در این حوزه چقدر قابل اتکاست، و برای اتکاپذیر کردنش چه چیزی
لازم است؟} این پرسش در آغاز کار مطرح نبود و از دل کار بیرون آمد.

\section{سهم این پایان‌نامه}

سهم اصلی این کار \textbf{روش‌شناختی} است، نه کارایی. این را باید در آغاز نوشت، چون خواننده‌ای
که به دنبال عددی بهتر آمده است باید بداند کجا را نگاه کند و کجا را نه.

\textbf{یکم، قاعدهٔ واسنجی همتاشده و سابقهٔ عمل آن.} هیچ اختلافی به بازنمایی نسبت داده نمی‌شود
مگر آن‌که هر دو مدل زیر نقاط برش برازش‌متقاطع سنجیده شوند. این قاعده تا امروز چهار ادعا را در
\emph{دو جهت} برگردانده و دو بار علامت اختلاف را وارونه کرده است، و سابقه‌اش در فصل
\ref{ch:results} می‌آید. مهم‌تر از هر تک‌نتیجه، این است که همان قاعده یک نتیجهٔ ظاهراً منفی —
کوری نسبت به درجهٔ خفیف — را به مسئله‌ای بی‌هزینه و حل‌شدنی تبدیل کرد.

\textbf{دوم، یک سقف نظارت اندازه‌گیری‌شده برای ادم ماکولا.} شش مداخلهٔ مستقل، از جمله
مداخله‌ای که مستقیماً ناحیهٔ تعریف‌کنندهٔ درجهٔ بالینی را جدا می‌کند، بازهٔ اطمینان را از صفر
دور نکردند. این نتیجه با شاهدی بیرونی تقویت می‌شود که همهٔ توضیح‌های جایگزین را می‌بندد: یک
پیکرهٔ ادم ماکولای ساخته‌شده برای همین کار، با \textbf{\num{۲۹۳۸}} تصویر، باز هم تنها
\textbf{\num{۹٫۵}} درصد نمونه در درجهٔ میانی دارد \cite{tang2026mmrdr}.

\textbf{سوم، تبار داده به‌عنوان سهم، نه به‌عنوان اعتراف.} شش اشکال تبار داده در این پژوهش هیچ
خطای زمان اجرا تولید نکردند و در سکوت عدد غلط ساختند. سازوکارهایی که پاسخ آن‌هاست در فصل
\ref{ch:method} می‌آید، و \textbf{یکی از پنج نمونهٔ آن فصل، شکست خودِ این پایان‌نامه است}.

\textbf{چهارم، یک نتیجهٔ منفی که راهی را می‌بندد.} درجهٔ ادم ماکولا با تعریف هندسی منتشرشده‌اش
از روی حاشیه‌نویسی‌های منتشرشده بازسازی نمی‌شود. این را فصل \ref{sec:parta} نشان می‌دهد و
پیامدش این است که مسیر ساختن دادهٔ سه‌کلاسهٔ بیشتر از روی ماسک‌های قطعه‌بندی بسته است.

\section{آنچه این پایان‌نامه ادعا نمی‌کند}

\textbf{هیچ عددی در این متن با کار منتشرشده قابل مقایسه نیست.} تقسیم این پژوهش اختصاصی و
منجمد است و بر پایهٔ خوشه‌بندی ادراکی ساخته شده تا نشت میان‌پیکره‌ای را ببندد؛ همین کار آن را
با هر تقسیم منتشرشده‌ای ناهم‌سنجه می‌کند. این محدودیت در فصل \ref{ch:results} و فصل
\ref{ch:discussion} تکرار می‌شود، چون تنها جایی که می‌تواند آسیب بزند همان جایی است که
خواننده عادت دارد عددها را کنار هم بگذارد.

\textbf{و ادعای برتری معماری هم در کار نیست.} از نه مداخلهٔ آزموده‌شده، هشت مداخله ابطال شدند.
آنچه می‌ماند رویه‌ای است برای رسیدن به عددی که بتوان از آن دفاع کرد.

\section{ساختار پایان‌نامه}

فصل \ref{ch:related} ادبیات را نه بر پایهٔ موضوع، که بر پایهٔ \emph{تناقض‌ها} سازمان می‌دهد:
سه جایی که نتیجه‌های منتشرشده با هم نمی‌خوانند، و این‌که هر تناقض در کدام خانهٔ آزموده‌نشده
حل می‌شود. فصل \ref{ch:method} پروتکل ارزیابی، داده، معماری و سازوکارهای تبار داده را می‌آورد؛
معماری آخرین چیزی است که در آن فصل اهمیت پیدا می‌کند. فصل \ref{ch:results} نتایج را می‌آورد،
با ترتیبی که از «خط لوله در برابر چه چیزی سنجیده می‌شود» آغاز می‌کند. فصل \ref{sec:parta}
گزارش مستقل نتیجهٔ منفی بازسازی درجه است. فصل \ref{ch:discussion} بحث، محدودیت‌ها و کارهای
آینده را جمع می‌بندد.
